\documentclass[10pt]{article}

\usepackage[utf8]{inputenc}

\usepackage[T1]{fontenc}

\usepackage{amsmath}

\usepackage{amsfonts}

\usepackage{amssymb}

\usepackage[version=4]{mhchem}

\usepackage{stmaryrd}

\usepackage{graphicx}

\usepackage[export]{adjustbox}

\graphicspath{ {./images/} }

\usepackage{bbold}



\begin{document}

Переход от одной пары параметров с ее потенциалом к другой паре с соответствующим потенциалом задается с помощью Лежандра преобразования. Так, при переходе от пары $(s, v)$ к паре $(T, v)$ потенциал $F(T, v)$ этой пары равен



\[

F(T, v)=E(s(T), v)-s(T) T,

\]



где $s(T)$ находится из уравнения (1), т. е. $F(T, v)$ с точностьо до знака совпадает с преобразованием Лежандра функции $E(s, v)$ как функции переменной $s$.



При содержательном построении термодинамики с помощью равновесных гиббсовских ансамблей Т. п. могут быть выражены с помощью термодинамич. предела, деленного на объем логарифма статистич. суммы (и его производных) какого-нибудь из гиббсовских ансамблей. Напр., свободная әнергия Гельмгольца равна



\[

F(T, v)=\lim _{N \rightarrow \infty} \frac{|\Lambda|}{N} \rightarrow v,

\]



где $Z(T, N, \Lambda)$ - статистическая сумма малого канонич. ансамбля для системы из $N$ частиц, заключенных в области $\Lambda(|\Lambda|-$ объем этой области $\Lambda)$, при фиксированном значении температуры $T$ (см. [3]).



Лum.: [1] Л а н д а у Л. Д., Л и ф шї и Е. М., Статистическая механика, 2 изд., М., 1964 (Теоретич. физика, т. 5),



\includegraphics[max width=\textwidth, center]{2023_07_09_755bff67b847c3b07e7fg-1(1)}

исчисление, М., $1961 ;$; 1 Р ю ә л ь Д., Статистическая механи-

ка. Строгие результаты, пер. с англ., м., 1971. P. A. Минлос. ТЕРМОДИНАМИЧЕСКИЙ ПРЕДЕЈЬНЫЙ ПЕРЕХОД - основной прием статистич. физики, заклютаюпийся в том, что для изучения большой (но конечной) физич. системы ее аппроксимируют нек-рой бесконечной идеализированной системой. Так, напр., систему, состоящую из $N$ частиц (молекул), заполняющих ограниченную область $\Lambda$ пространства $\mathbb{R}^{3}$, при большом $N$ и большой (сравнительно с размерами молекул) области $\Lambda$ заменяют системой из бесконечного числа тех же молекул, заполняющих все пространство, так что при әтом свойства и характеристики конечной системы (характер динамики, свойства равновесных ансамблей и т. д.) асимптотически (по $N$ и по $\Lambda$ ) близки к аналогичным свойствам и характеристикам предельной системы.



Лит.: [1] Р ю ә л ь Д., Статистическая механика. Строгие

М. результаты, пер. с англ., М., 1970; [2] Гиббсовские состояния

в статистической физик, пер. с англ., М., 1978; [3] М и нл о с Р. А., «Успехи матем. наук», 1968, т. 23 , в. 1, с. 133-90.



\includegraphics[max width=\textwidth, center]{2023_07_09_755bff67b847c3b07e7fg-1(3)}

нельзя представить в виде пересечения строго больших чем $I$ правого частного $r(I, A)$ и идеала $B$. Все неприводимые идеалы терциарны. В нётеровых кольцах терциарность совпадает с примарностью (см. Аддитивная теория идеалов, Примарный идеал, Примарное разложение).



Пусть кольцо $R$ удовлетворяет условию максимальности для левых и правых частных идеалов и каждый идеал разложим в пересечение конечного числа неприводимых идеалов. Тогда для каждого идеала $Q$ существует т е р ц и а р ный ра ди к а л ter $(Q)-$ наибольший среди таких идеалов $T$ в $R$, что для лобого идеала $B$



\[

r(Q, T) \cap B=Q \Rightarrow B=Q .

\]



Как и для примарных идеалов, для Т. и. сшраведливы теорема пересечения, теорема существования и теорема единственности.



Анализ свойств правых и левых частных (идеалов кольца, подмодулей модуля и др.) приводят к системам с частными, в к-рых естественным образом вводится общее понятие $S$-примарности и $S$-примарного радикала. Это позволяет сформулировать в качестве аксиом «теоремы пересечения», «существования» и «единственности». При таком подходе терциарность является единственной примарностью, для к-рой справедливы все эти три теоремы, т. е. является единственным «хорошим» обобщением классич. примарности (см. [1], [2]).



\begin{center}

\includegraphics[max width=\textwidth]{2023_07_09_755bff67b847c3b07e7fg-1}

\end{center}



\begin{center}

\includegraphics[max width=\textwidth]{2023_07_09_755bff67b847c3b07e7fg-1(2)}

\end{center}



TECT в к и б е р н е т и к е- одно из важнейших средств логич. анализа информации. Аппарат T. первоначально был использован в задаче контроля работы электрич. схем. Позже на основе Т. были разработаны эффективные алгоритмы распознавания образов. Тестовый подход успешно применяется во многих областях математики. Основные задачи с Т. могут быть сформулированы следующим образом.



\begin{enumerate}

  \item Дана прямоугольная таблица, содержацая $s$ строк и $m$ столбцов. Строки характеризуются вопросами (признаками) $e_{1}$, $\ldots, e_{s}$ из нек-рого множества $E$,столбцы - объектами (образами) $f_{1}, \ldots \dot{\text {, }}$

$f_{m}$ из множества $F$.

\end{enumerate}



\includegraphics[max width=\textwidth, center]{2023_07_09_755bff67b847c3b07e7fg-1(4)}

В клетке, лежащей на пересечении $i$-й строки и $j$-го столбда, находится ответ $f_{j}\left(e_{i}\right)$ (значение $i$-го признака для $j$-го образа), принадлежащий нек-рому множеству $G$. Таким образом, поскольку $e_{i_{1}} \neq e_{i_{2}}$ при $i_{1} \neq i_{2}$, можно рассматривать $j$-й столбец $(1 \leqslant j \leqslant m)$ как столбец, задающий функцию $f_{j}(x)$. Естественно считать, что столбцы таблицы попарно различны. Так как природа множеств $E$ и $G$ существенного значения не имеет, то в дальнейшем предполагается $E=\{0,1, \ldots, s-1\}$ и $G=\{0,1, \ldots, k-1\}$. В нек-рых случаях на множествах $E$ и $G$ задается частичный порядок $\leqslant$. Иногда столбцам приписываются вероятности $p_{1}, \ldots . p_{m}\left(\Sigma_{i} p_{i}=1\right)$.



\begin{enumerate}

  \setcounter{enumi}{1}

  \item Задается цель логич. анализа таблицы. Для этого фиксируется нек-рое подмножество $\mathfrak{A}$ пар $(i, j), i \neq j$, номеров столбцов. В частности, если множество $F$ разбито на классы, то $(i, j) \in \mathfrak{\Re}$ тогда и только тогда, когда $f_{i}$ и $f_{j}$ принадлежат одному классу, $i \neq j$. Подмножество $9 \mathcal{2}$ можно интерпретировать как отношение или как нек-рое свойство.

\end{enumerate}



Имеется два специальных случая: 1) $\mathfrak{R}=\{(1, j)\}$, $j=2, \ldots, m$ и $f_{1}$ - выделенная функция (этот случай связан с задачей о проверке), 2) $\mathfrak{l}=\{(i, j)\}, i, j=1$, ..., $m, i \neq j$, относится к задаче о диагностике (нолного различения).



Цель логич. анализа можно сформулировать следующим образом: указать процедуру, к-рая для любых $i$ и $j$ таких, что $(i, j) \in \mathfrak{N}$, позволяла бы отличить образ $f_{i}$ от образа $f_{j}$. Если $\mathfrak{l}$ соответствует разбшению $F$ на классы, то задача эквивалентна задаче об отнесении произвольно заданной функции $f$ из $F$ соответствующему классу. Ниже в основном говорится об этой задаче. Сформулированный вопрос решается тривиально, если полностью известна таблица и все известно о функции $f$. В реальных задачах получение информации о функции $f$ возможно, но при определенных затратах. К роме того, для обширного класса задач вся таблица либо не известна, либо настолько велика, что мы не в состоянии с ней оперировать. Поәтому ограничиваются ее нек-рым фрагментом - т. н. п р е д с т а в и т е л ь н о й в ыб о р о й, и задачу ставят в әвристич. варианте.



\begin{enumerate}

  \setcounter{enumi}{2}

  \item Указываются допустимые средства решения задачи. Пусть загадана нек-рая функция $f$ из $F$. Требуется путем задавания вопросов, называя нек-рые строки $e_{i_{1}}, \ldots, e_{i_{r}}$, по ответам $f\left(e_{i_{1}}\right), \ldots, f\left(e_{i_{r}}\right)$ узнать, к какому классу принадлежит $f$. Данный опрос может осу ществляться либо при помощи т. н. бе з у сло в н ого ә к с пе р и м е н т а (см. Эксперимеитьл с авто-

\end{enumerate}



\end{document}